% !TEX TS-program = pdflatexmk
\documentclass[12pt]{amsart}

%\usepackage[parfill]{parskip}    % Activate to begin paragraphs with an empty line rather than an indent

\usepackage[margin=1in]{geometry}

\usepackage{amsmath,amssymb,amsthm,latexsym,graphicx}
\usepackage[normalem]{ulem}
\usepackage{setspace} %used for doublespacing, etc.
\usepackage{hyperref}
\usepackage[shortlabels]{enumitem}
\usepackage{cancel}
\usepackage[dvipsnames,usenames]{color}
\usepackage[all]{xy}
\DeclareGraphicsRule{.tif}{png}{.png}{`convert #1 `dirname #1`/`basename #1 .tif`.png}

%Some useful environments.
\newtheorem{theorem}{Theorem}
\newtheorem{corollary}[theorem]{Corollary}
\newtheorem{conjecture}[theorem]{Conjecture}
\newtheorem{lemma}[theorem]{Lemma}
\newtheorem{proposition}[theorem]{Proposition}
\newtheorem{definition}[theorem]{Definition}
\newtheorem{example}[theorem]{Example}
\newtheorem{axiom}{Axiom}
\theoremstyle{remark}
\newtheorem{remark}{Remark}
\newtheorem*{exercise}{Exercise}%[section]

%Some useful shortcuts for our favorite fields
\newcommand{\RR}{{\mathbb R}}
\newcommand{\NN}{{\mathbb N}}
\newcommand{\ZZ}{{\mathbb Z}}
\newcommand{\QQ}{{\mathbb Q}}
\newcommand{\CC}{{\mathbb C}}

%Some useful shortcuts for formatting lists
\newcommand{\bc}{\begin{center}}
\newcommand{\ec}{\end{center}}
\newcommand{\be}{\begin{enumerate}}
\newcommand{\ee}{\end{enumerate}}
\newcommand{\bi}{\begin{itemize}}
\newcommand{\ei}{\end{itemize}}

%Some useful shortcuts for formatting mathematical symbols
\newcommand{\ol}[1]{\overline{#1}}
\newcommand{\oimp}[1]{\overset{#1}{\Longleftrightarrow}} %labeled iff symbol
\newcommand{\bv}[1]{\ensuremath{ \vec{\mathbf{#1}}} } %makes a vector.
\newcommand{\mc}[1]{\ensuremath{\mathcal{#1}}} %put something in caligraphic font
\newcommand{\mpg}[1]{\marginpar{ #1}} %to put comments in margins
\newcommand{\bsl}[1]{\texttt{\symbol{92}{\em #1}}} %for backslashes.
\newcommand{\normale}{\trianglelefteq}
\newcommand{\normal}{\triangleleft}

%Commenting tools
\usepackage{soul}
\definecolor{highlight}{rgb}{1,0.6,0.6}
\sethlcolor{highlight}
\newcommand{\hlm}[1]{\colorbox{highlight}{$\displaystyle #1$}}
\newtheoremstyle{mycomment}{\topsep}{-0in}{\small \itshape \sffamily}{}{\small \itshape\sffamily}{:}{.5em}{}
\theoremstyle{mycomment}
\newtheorem*{acomment}{\color{BrickRed}{Comment}}
\newcommand{\com}[1]{{\color{OliveGreen}\begin{acomment}{#1} %#2 \color{black} 
\end{acomment}\noindent}}
%\newcommand{\com}[1]{{\color{BrickRed}{\\ Comment:}\color{OliveGreen}{#1} \\}}
\newcommand{\red}[1]{{\color{BrickRed} #1}}
\newcommand{\blue}[1]{{\color{MidnightBlue}#1}}
\newcommand{\green}[1]{{\color{OliveGreen}#1}}
\newcommand{\mwrong}[2]{\red{\cancel{#1}}\green{#2}}
\newcommand{\wrong}[2]{\red{\sout{#1}}\green{#2}}
\definecolor{OliveGreen}{rgb}{.3,.5,.2}
\definecolor{MidnightBlue}{rgb}{.3,.4,.6}

\title{Assignment \#8 10/27/22}

\author{Coordinator: Stefano Fochesatto}
\begin{document}
\maketitle
%\setstretch{2.5} %Use for 2.5 spacing
\doublespacing %Use for double spacing

\section*{Section 5.1}
\begin{exercise}[5.1.10 --- Noah Palmer] Let $p$ be a prime. Let $A$ and $B$ be two cyclic groups of order $p$ with generators $x$ and $y$, respectively. Set $E=A\times B$ so that $E$ is the elementary abelian group of order $p^2: E_{p^2}$. Prove that the distinct subgroups of $E$ of order $p$ are
$$\langle x\rangle,\hspace{0.5cm}\langle xy\rangle,\hspace{0.5cm}\ldots \hspace{0.5cm},\langle xy^{p-1}\rangle,\hspace{0.5cm}\langle y\rangle$$
\begin{proof}
Let $p$ be a prime, $\langle x\rangle=A$ and $\langle y\rangle=B$ be two cyclic groups of order $p$, and set $E=A\times B$ so that $E$ is the elementary abelian group of order $p^2$. Since $A$ and $B$ are cyclic groups of prime order it follows that each element of $A$ and $B$ have order $p$. Hence, each of the subgroups
$$\langle x\rangle,\hspace{0.5cm}\langle xy\rangle,\hspace{0.5cm}\ldots \hspace{0.5cm},\langle xy^{p-1}\rangle,\hspace{0.5cm}\langle y\rangle$$
will have order $p$. Additionally, each of these groups are unique since each $y^{i}$ is unique within its own cyclic group which implies that the subgroups generated by $\langle(1,1,\ldots,x,\ldots, y^i,\ldots 1)\rangle$ will be unique for each $1\leq i<p$. From example $3$ in section $5.1$ of Dummit and Foote the number of subgroups of order $p$ in an elemntary abelian group of order $p^2$ is $p+1$. Since, we have given a list of exactly $p+1$ distinct groups, we're done. 

\end{proof}
\end{exercise}

\section*{Section 5.2}
\begin{exercise}[5.2.4 --- Noah Palmer] In each of parts $(a)$ to $(d)$ determine which pairs of abelian groups listed are isomorphic.
\begin{enumerate}[(a)]
\item $\{4,9\}$, $\{6,6\}$, $\{8,3\}$, $\{9,4\}$, $\{6,4\}$, $\{64\}$.
\item $\{2^2, 2\cdot3^2\}$, $\{2^2\cdot3, 2\cdot3\}$, $\{2^3\cdot3^2\}$, $\{2^2\cdot3^2,2\}$.
\item $\{5^2\cdot7^2,3^2\cdot5\cdot7\}$, $\{3^2\cdot5^2\cdot7,5\cdot7^2\}$, $\{3\cdot5^2,7^2,3\cdot5\cdot7\}$, $\{5^2\cdot7,3^2\cdot5,7^2\}$.
\item $\{2^2\cdot5\cdot7, 2^3\cdot5^3,2\cdot5^2\}$, $\{2^3\cdot5^3\cdot7,2^3\cdot5^3\}$, $\{2^2,2\cdot7,2^3,5^3,5^3\}$, $\{2\cdot5^3,2^2\cdot5^3,2^3,7\}$.
\end{enumerate}
\begin{proof}
We use the fact that $Z_m\times Z_n\cong Z_{mn}$ if and only if $\gcd(m,n)=1$ to rearrange the abelian groups into ismorphic pairs, where possible.
\begin{enumerate}[(a)]
\item Note that $ Z_4\times Z_9\cong Z_9\times Z_4$. Thus, $\{4,9\}\cong\{9,4\}$.
\item Note that $ Z_{2^2}\times Z_{2\cdot3^2}\cong Z_{2^2\cdot3^2}\times Z_2$. Thus, $\{2^2, 2\cdot3^2\}\cong\{2^2\cdot3^2,2\}$.
\item Note that $Z_{5^2\cdot 7^2}\times Z_{3^2\cdot 5\cdot 7}\cong Z_{3^2\cdot 5^2\cdot 7}\times Z_{5\cdot 7^2}$. Also, $Z_{3^2\cdot 5^2\cdot 7}\times Z_{5\cdot 7^2}\cong Z_{5^2\cdot 7}\times Z_{3^2\cdot 5}\times Z_{7^2}$. Thus, $\{5^2\cdot7^2,3^2\cdot5\cdot7\}\cong\{3^2\cdot5^2\cdot7,5\cdot7^2\}\cong\{5^2\cdot7,3^2\cdot5,7^2\}$.
\item Note that $ Z_{2^2}\times Z_{2\cdot7}\times Z_{2^3}\times Z_{5^3}\times Z_{5^3}\cong Z_{2^2\cdot5^3}\times Z_{2\cdot7\cdot5^3}\times Z_{2^3}\cong Z_{2\cdot5^3}\times Z_{2^2\cdot5^3}\times Z_{2^3}\times Z_7$. Therefore, $\{2^2,2\cdot7,2^3,5^3,5^3\}\cong\{2\cdot5^3,2^2\cdot5^3,2^3,7\}$.
\end{enumerate}
\end{proof}
\end{exercise}





\begin{exercise}[5.2.9 --- Noah Palmer] Let $A=Z_{60}\times Z_{45}\times Z_{12}\times Z_{36}$. Find the number of elements of order $2$ and the number of subgroups of index $2$ in $A$.
\begin{proof}
Begin by rewriting the direct product in its invariant form $Z_{60}\times Z_{45}\times Z_{12}\times Z_{36}\cong Z_{180}\times Z_{60}\times Z_3\times Z_{36}\cong Z_{180}\times Z_{180}\times Z_{12}\times Z_{3}$. Now, note that each of $Z_{180}$ and $Z_{12}$ has an element of order $2$ whereas $Z_{3}$ does not. Hence, there will be $2^3$ elements of order less than or equal to $2$. From these we must exclude the identity since it has order $1$. Thus, there are $7$ elements of order $2$.

To find the subgroups of index $2$ in $A$ consider the elementary divisor factorization of $A=Z_{5}\times Z_{5}\times Z_{4}\times Z_{4}\times Z_4\times Z_{9}\times Z_{9}\times Z_3\times Z_3$. Now, any subgroup of order $583,200$ will have index $2$ in $A$. From the above factorization, and the prime factorization of $583,200$, these subgroups will have the form $A^\prime\times Z_{5}\times Z_{5}\times Z_{9}\times Z_{9}\times Z_3\times Z_3$ where $A^\prime$ is a subgroup of $Z_{4}\times Z_{4}\times Z_4$ with order $32$. There are three ways to select such a subgroup (e.g. $Z_4\times Z_4\times Z_2$, $Z_2\times Z_4\times Z_4$, and $Z_4\times Z_2\times Z_4$). Thus, there are three subgroups of index $2$ in $A$.
\begin{center}

\end{center}
\end{proof}
\end{exercise}
 \end{document}
 \end

  